\documentclass[11pt]{article}
\usepackage[T1]{fontenc}
\usepackage[utf8]{inputenc}
\usepackage{lmodern}
\usepackage{microtype}
\usepackage[margin=0.72in,headheight=15pt]{geometry}
\usepackage[table]{xcolor}
\usepackage{array,booktabs,longtable,tabularx}
\usepackage{amssymb}
\usepackage{enumitem}
\usepackage[most]{tcolorbox}
\usepackage{fancyhdr}
\usepackage{titlesec}
\usepackage{needspace}
\usepackage{ragged2e}
\usepackage{hyperref}
\usepackage{bookmark}
\usepackage{lastpage}

\definecolor{Navy}{HTML}{17324D}
\definecolor{Teal}{HTML}{157A7A}
\definecolor{CueGray}{HTML}{EEF2F4}
\definecolor{SoftBlue}{HTML}{E9F3F8}
\definecolor{SoftGreen}{HTML}{EAF5F0}
\definecolor{SoftAmber}{HTML}{FFF4D6}
\definecolor{RuleGray}{HTML}{B7C3CA}
\definecolor{TextGray}{HTML}{263238}

\hypersetup{colorlinks=true,linkcolor=Navy,urlcolor=Teal,citecolor=Navy,pdfborder={0 0 0}}
\color{TextGray}
\setlength{\parindent}{0pt}
\setlength{\parskip}{4pt}
\setlength{\emergencystretch}{3em}
\hfuzz=0.2pt
\hbadness=2000
\setlist{nosep,leftmargin=1.4em}
\setlength{\LTpre}{4pt}
\setlength{\LTpost}{6pt}
\renewcommand{\arraystretch}{1.2}

\titleformat{\section}{\Large\bfseries\color{Navy}\RaggedRight}{\thesection}{0.55em}{}
\titleformat{\subsection}{\normalsize\bfseries\color{Teal}\RaggedRight}{\thesubsection}{0.5em}{}
\titlespacing*{\section}{0pt}{1.3ex plus .3ex}{.7ex}
\titlespacing*{\subsection}{0pt}{1.1ex plus .2ex}{.5ex}

\pagestyle{fancy}
\fancyhf{}
\lhead{\small\color{Navy}\textbf{Cybersecurity Cornell Notes}}
\rhead{\small\color{Teal}\ChapterMark}
\cfoot{\small\thepage\ of \pageref*{LastPage}}
\renewcommand{\headrulewidth}{0.5pt}

\newtcolorbox{orientationbox}{enhanced,breakable,colback=SoftBlue,colframe=Teal,boxrule=.7pt,arc=2mm,left=2mm,right=2mm,top=1.5mm,bottom=1.5mm}
\newtcolorbox{topicsummary}[1][]{enhanced,breakable,colback=CueGray,colframe=RuleGray,title={Topic Summary},fonttitle=\bfseries\color{Navy},boxrule=.5pt,arc=1.5mm,left=2mm,right=2mm,top=1mm,bottom=1mm,#1}
\newtcolorbox{defensebox}{enhanced,breakable,colback=SoftGreen,colframe=Teal,title={Defensive Guidance},fonttitle=\bfseries,boxrule=.7pt,arc=2mm}
\newtcolorbox{historybox}{enhanced,breakable,colback=SoftAmber,colframe=orange!70!black,title={Historical Context},fonttitle=\bfseries,boxrule=.7pt,arc=2mm}
\newtcolorbox{synthesisbox}{enhanced,breakable,colback=Navy!5,colframe=Navy,title={Chapter Synthesis},fonttitle=\bfseries,boxrule=.8pt,arc=2mm}
\newtcolorbox{takeawaybox}{enhanced,colback=Teal!8,colframe=Teal,title={One-Sentence Takeaway},fonttitle=\bfseries,boxrule=.9pt,arc=2mm}

\newcommand{\CornellHeader}{%
\rowcolor{Navy}\color{white}\bfseries Cue / Recall Prompt & \color{white}\bfseries Notes / Analysis\\\hline}
\newcommand{\CornellRow}[2]{%
\cellcolor{CueGray}\RaggedRight\textbf{#1} & \RaggedRight #2\\\hline}

\newcommand{\ChapterMark}{Chapter 71}
\title{\textcolor{Navy}{Chapter 71: Physical Security}\\[2mm]\large Cornell Notes}
\author{}
\date{}
\hypersetup{pdftitle={Chapter 71: Physical Security - Cornell Notes},pdfauthor={},pdfsubject={Cybersecurity study notes},pdfkeywords={physical, access, threats, and equipment}}
\begin{document}
\maketitle
\thispagestyle{fancy}
\vspace{-1.5em}
\begin{orientationbox}
\textbf{Chapter orientation.} This chapter examines physical security within the broader domain of cyber-physical security. Its central problem is how to reason about physical, access, threats, and equipment while keeping security objectives, operating assumptions, evidence, and recovery connected. A practitioner should approach the material through physical consequences, safety, availability, environmental dependencies, remote connectivity, maintenance, and recovery. Useful prerequisites include basic risk, threat, control, and systems concepts.
\end{orientationbox}
\section*{Learning Objectives}
\begin{itemize}
\item Explain the security purpose and scope of Physical Security.
\item Identify the assets, actors, assumptions, and trust boundaries associated with physical, access, and threats.
\item Distinguish Chapter orientation from Overview and explain where their responsibilities intersect.
\item Analyze how failures in physical, access, and threats can affect confidentiality, integrity, availability, privacy, or safety.
\item Evaluate relevant controls through physical consequences, safety, availability, environmental dependencies, remote connectivity, maintenance, and recovery.
\item Audit an implementation using asset and dependency maps, physical-access records, sensor telemetry, safety constraints, maintenance logs, and recovery exercises.
\item Prioritize remediation according to exploitability, consequence, control strength, and recovery readiness.
\item Verify that corrective actions change observable risk rather than documentation alone.
\end{itemize}
\section*{Cornell Cue-and-Notes}
\Needspace{8\baselineskip}
\subsection*{Chapter orientation}
\begin{longtable}{>{\RaggedRight\arraybackslash}p{0.29\textwidth}|>{\RaggedRight\arraybackslash}p{0.65\textwidth}}
\CornellHeader
\CornellRow{What security problem does Chapter orientation address?}{\textbf{Core idea.} Treat Chapter orientation as a structured part of Physical Security, with particular attention to physical, theft, protects, and premises. \textbf{Analytical lens.} This topic establishes a component of the chapter's argument; connect its definitions and mechanisms to a testable security objective. For physical, theft, and protects, model safety and physical consequences alongside conventional information-security effects. \textbf{Security reasoning.} Identify what must remain protected, who or what can alter the expected state, which assumptions cross a trust boundary, and how failure changes confidentiality, integrity, availability, privacy, or safety. \textbf{Application.} The practitioner should evaluate cyber and physical failure together, enforce safe states, and test operational recovery. \textbf{Evidence.} Seek asset and dependency maps, physical-access records, sensor telemetry, safety constraints, maintenance logs, and recovery exercises.}
\end{longtable}
\begin{topicsummary}
Chapter orientation is best remembered as the connection between physical, theft, protects, and premises and the chapter's larger control objective. A defensible review states the assumption, expected behavior, evidence, failure consequence, and required response.
\end{topicsummary}
\Needspace{8\baselineskip}
\subsection*{Overview}
\begin{longtable}{>{\RaggedRight\arraybackslash}p{0.29\textwidth}|>{\RaggedRight\arraybackslash}p{0.65\textwidth}}
\CornellHeader
\CornellRow{Which assumptions matter in Overview?}{\textbf{Core idea.} Treat Overview as a structured part of Physical Security, with particular attention to physical, storage, assets, and role. \textbf{Analytical lens.} This topic is architecture-centered: locate components, interfaces, trust boundaries, dependencies, control points, and failure propagation. For physical, storage, and assets, model safety and physical consequences alongside conventional information-security effects. \textbf{Security reasoning.} Identify what must remain protected, who or what can alter the expected state, which assumptions cross a trust boundary, and how failure changes confidentiality, integrity, availability, privacy, or safety. \textbf{Application.} The practitioner should evaluate cyber and physical failure together, enforce safe states, and test operational recovery. \textbf{Evidence.} Seek asset and dependency maps, physical-access records, sensor telemetry, safety constraints, maintenance logs, and recovery exercises.}
\end{longtable}
\begin{topicsummary}
Overview is best remembered as the connection between physical, storage, assets, and role and the chapter's larger control objective. A defensible review states the assumption, expected behavior, evidence, failure consequence, and required response.
\end{topicsummary}
\Needspace{8\baselineskip}
\subsection*{Physical Security Threats}
\begin{longtable}{>{\RaggedRight\arraybackslash}p{0.29\textwidth}|>{\RaggedRight\arraybackslash}p{0.65\textwidth}}
\CornellHeader
\CornellRow{How should a reviewer evaluate Physical Security Threats?}{\textbf{Core idea.} Treat Physical Security Threats as a structured part of Physical Security, with particular attention to damage, threats, equipment, and temperature. \textbf{Analytical lens.} This topic is threat-centered: separate actor capability, reachable attack surface, prerequisites, execution path, and consequence. For damage, threats, and equipment, model safety and physical consequences alongside conventional information-security effects. \textbf{Security reasoning.} Identify what must remain protected, who or what can alter the expected state, which assumptions cross a trust boundary, and how failure changes confidentiality, integrity, availability, privacy, or safety. \textbf{Application.} The practitioner should evaluate cyber and physical failure together, enforce safe states, and test operational recovery. \textbf{Evidence.} Seek asset and dependency maps, physical-access records, sensor telemetry, safety constraints, maintenance logs, and recovery exercises. Related subtopics include Natural Disasters, Environmental Threats, Inappropriate Temperature and Humidity, and Fire and Smoke.}
\end{longtable}
\begin{topicsummary}
Physical Security Threats is best remembered as the connection between damage, threats, equipment, and temperature and the chapter's larger control objective. A defensible review states the assumption, expected behavior, evidence, failure consequence, and required response.
\end{topicsummary}
\Needspace{8\baselineskip}
\subsection*{Physical Security Prevention And Mitigation Measures}
\begin{longtable}{>{\RaggedRight\arraybackslash}p{0.29\textwidth}|>{\RaggedRight\arraybackslash}p{0.65\textwidth}}
\CornellHeader
\CornellRow{Why can Physical Security Prevention And Mitigation Measures fail in practice?}{\textbf{Core idea.} Treat Physical Security Prevention And Mitigation Measures as a structured part of Physical Security, with particular attention to power, equipment, threats, and supply. \textbf{Analytical lens.} This topic is control-centered: map each safeguard to a specific failure mode, then test independence, coverage, and bypass paths. For power, equipment, and threats, model safety and physical consequences alongside conventional information-security effects. \textbf{Security reasoning.} Identify what must remain protected, who or what can alter the expected state, which assumptions cross a trust boundary, and how failure changes confidentiality, integrity, availability, privacy, or safety. \textbf{Application.} The practitioner should evaluate cyber and physical failure together, enforce safe states, and test operational recovery. \textbf{Evidence.} Seek asset and dependency maps, physical-access records, sensor telemetry, safety constraints, maintenance logs, and recovery exercises. Related subtopics include Electromagnetic Interference, Environmental Threats, Inappropriate Temperature and Humidity, and Water Damage.}
\end{longtable}
\begin{topicsummary}
Physical Security Prevention And Mitigation Measures is best remembered as the connection between power, equipment, threats, and supply and the chapter's larger control objective. A defensible review states the assumption, expected behavior, evidence, failure consequence, and required response.
\end{topicsummary}
\Needspace{8\baselineskip}
\subsection*{Recovery From Physical Security Breaches}
\begin{longtable}{>{\RaggedRight\arraybackslash}p{0.29\textwidth}|>{\RaggedRight\arraybackslash}p{0.65\textwidth}}
\CornellHeader
\CornellRow{What security problem does Recovery From Physical Security Breaches address?}{\textbf{Core idea.} Treat Recovery From Physical Security Breaches as a structured part of Physical Security, with particular attention to site, recovery, physical, and nature. \textbf{Analytical lens.} This topic is resilience-centered: identify failure domains, acceptable degradation, recovery objectives, dependencies, and tested restoration paths. For site, recovery, and physical, model safety and physical consequences alongside conventional information-security effects. \textbf{Security reasoning.} Identify what must remain protected, who or what can alter the expected state, which assumptions cross a trust boundary, and how failure changes confidentiality, integrity, availability, privacy, or safety. \textbf{Application.} The practitioner should evaluate cyber and physical failure together, enforce safe states, and test operational recovery. \textbf{Evidence.} Seek asset and dependency maps, physical-access records, sensor telemetry, safety constraints, maintenance logs, and recovery exercises. Related subtopics include Human-Caused Physical Threats.}
\end{longtable}
\begin{topicsummary}
Recovery From Physical Security Breaches is best remembered as the connection between site, recovery, physical, and nature and the chapter's larger control objective. A defensible review states the assumption, expected behavior, evidence, failure consequence, and required response.
\end{topicsummary}
\Needspace{8\baselineskip}
\subsection*{Threat Assessment, Planning, And Plan Implementation}
\begin{longtable}{>{\RaggedRight\arraybackslash}p{0.29\textwidth}|>{\RaggedRight\arraybackslash}p{0.65\textwidth}}
\CornellHeader
\CornellRow{Which assumptions matter in Threat Assessment, Planning, And Plan Implementation?}{\textbf{Core idea.} Treat Threat Assessment, Planning, And Plan Implementation as a structured part of Physical Security, with particular attention to physical, access, resource, and plan. \textbf{Analytical lens.} This topic is threat-centered: separate actor capability, reachable attack surface, prerequisites, execution path, and consequence. For physical, access, and resource, model safety and physical consequences alongside conventional information-security effects. \textbf{Security reasoning.} Identify what must remain protected, who or what can alter the expected state, which assumptions cross a trust boundary, and how failure changes confidentiality, integrity, availability, privacy, or safety. \textbf{Application.} The practitioner should evaluate cyber and physical failure together, enforce safe states, and test operational recovery. \textbf{Evidence.} Seek asset and dependency maps, physical-access records, sensor telemetry, safety constraints, maintenance logs, and recovery exercises. Related subtopics include Threat Assessment.}
\end{longtable}
\begin{topicsummary}
Threat Assessment, Planning, And Plan Implementation is best remembered as the connection between physical, access, resource, and plan and the chapter's larger control objective. A defensible review states the assumption, expected behavior, evidence, failure consequence, and required response.
\end{topicsummary}
\Needspace{8\baselineskip}
\subsection*{Example: A Corporate Physical Security Policy}
\begin{longtable}{>{\RaggedRight\arraybackslash}p{0.29\textwidth}|>{\RaggedRight\arraybackslash}p{0.65\textwidth}}
\CornellHeader
\CornellRow{How should a reviewer evaluate Example: A Corporate Physical Security Policy?}{\textbf{Core idea.} Treat Example: A Corporate Physical Security Policy as a structured part of Physical Security, with particular attention to threats, company, costs, and physical. \textbf{Analytical lens.} This topic is governance-centered: connect required intent to accountable ownership, implementation, evidence, exceptions, and corrective action. For threats, company, and costs, model safety and physical consequences alongside conventional information-security effects. \textbf{Security reasoning.} Identify what must remain protected, who or what can alter the expected state, which assumptions cross a trust boundary, and how failure changes confidentiality, integrity, availability, privacy, or safety. \textbf{Application.} The practitioner should evaluate cyber and physical failure together, enforce safe states, and test operational recovery. \textbf{Evidence.} Seek asset and dependency maps, physical-access records, sensor telemetry, safety constraints, maintenance logs, and recovery exercises.}
\end{longtable}
\begin{topicsummary}
Example: A Corporate Physical Security Policy is best remembered as the connection between threats, company, costs, and physical and the chapter's larger control objective. A defensible review states the assumption, expected behavior, evidence, failure consequence, and required response.
\end{topicsummary}
\Needspace{8\baselineskip}
\subsection*{Integration Of Physical And Logical Security}
\begin{longtable}{>{\RaggedRight\arraybackslash}p{0.29\textwidth}|>{\RaggedRight\arraybackslash}p{0.65\textwidth}}
\CornellHeader
\CornellRow{Why can Integration Of Physical And Logical Security fail in practice?}{\textbf{Core idea.} Treat Integration Of Physical And Logical Security as a structured part of Physical Security, with particular attention to access, card, control, and physical. \textbf{Analytical lens.} This topic establishes a component of the chapter's argument; connect its definitions and mechanisms to a testable security objective. For access, card, and control, model safety and physical consequences alongside conventional information-security effects. \textbf{Security reasoning.} Identify what must remain protected, who or what can alter the expected state, which assumptions cross a trust boundary, and how failure changes confidentiality, integrity, availability, privacy, or safety. \textbf{Application.} The practitioner should evaluate cyber and physical failure together, enforce safe states, and test operational recovery. \textbf{Evidence.} Seek asset and dependency maps, physical-access records, sensor telemetry, safety constraints, maintenance logs, and recovery exercises. Related subtopics include Planning and Implementation.}
\end{longtable}
\begin{topicsummary}
Integration Of Physical And Logical Security is best remembered as the connection between access, card, control, and physical and the chapter's larger control objective. A defensible review states the assumption, expected behavior, evidence, failure consequence, and required response.
\end{topicsummary}
\Needspace{8\baselineskip}
\subsection*{Physical Security Checklist}
\begin{longtable}{>{\RaggedRight\arraybackslash}p{0.29\textwidth}|>{\RaggedRight\arraybackslash}p{0.65\textwidth}}
\CornellHeader
\CornellRow{What security problem does Physical Security Checklist address?}{\textbf{Core idea.} Treat Physical Security Checklist as a structured part of Physical Security, with particular attention to true false, physical, checklist, and natural. \textbf{Analytical lens.} This topic establishes a component of the chapter's argument; connect its definitions and mechanisms to a testable security objective. For true false, physical, and checklist, model safety and physical consequences alongside conventional information-security effects. \textbf{Security reasoning.} Identify what must remain protected, who or what can alter the expected state, which assumptions cross a trust boundary, and how failure changes confidentiality, integrity, availability, privacy, or safety. \textbf{Application.} The practitioner should evaluate cyber and physical failure together, enforce safe states, and test operational recovery. \textbf{Evidence.} Seek asset and dependency maps, physical-access records, sensor telemetry, safety constraints, maintenance logs, and recovery exercises.}
\end{longtable}
\begin{topicsummary}
Physical Security Checklist is best remembered as the connection between true false, physical, checklist, and natural and the chapter's larger control objective. A defensible review states the assumption, expected behavior, evidence, failure consequence, and required response.
\end{topicsummary}
\begin{historybox}
Some products, protocols, examples, threat observations, or legal references in this chapter are historically bounded. Retain the underlying threat model and control objective, but verify present applicability before treating a named implementation as current guidance.
\end{historybox}
\begin{defensebox}
Apply the chapter by making assets, authority, trust, expected behavior, and failure response explicit. In practice, evaluate cyber and physical failure together, enforce safe states, and test operational recovery. A control is not fully supported until its operation, monitoring, ownership, and recovery behavior are demonstrated with evidence.
\end{defensebox}
\section*{Threat-and-Control Map}
\begingroup\scriptsize\setlength{\tabcolsep}{2pt}
\begin{longtable}{>{\RaggedRight\arraybackslash}p{0.135\textwidth}>{\RaggedRight\arraybackslash}p{0.145\textwidth}>{\RaggedRight\arraybackslash}p{0.155\textwidth}>{\RaggedRight\arraybackslash}p{0.145\textwidth}>{\RaggedRight\arraybackslash}p{0.16\textwidth}>{\RaggedRight\arraybackslash}p{0.17\textwidth}}
\toprule\rowcolor{Navy}\color{white}\textbf{Asset} & \color{white}\textbf{Threat / Failure} & \color{white}\textbf{Vulnerability} & \color{white}\textbf{Impact} & \color{white}\textbf{Detection / Evidence} & \color{white}\textbf{Control}\\\midrule
\endfirsthead
\toprule\rowcolor{Navy}\color{white}\textbf{Asset} & \color{white}\textbf{Threat / Failure} & \color{white}\textbf{Vulnerability} & \color{white}\textbf{Impact} & \color{white}\textbf{Detection / Evidence} & \color{white}\textbf{Control}\\\midrule
\endhead
People, physical processes, connected devices, and essential services & Tampering, unsafe command, service disruption, surveillance, or cascading failure & Insecure remote access, fragile dependencies, weak update paths, or insufficient safety separation & Safety events, prolonged outage, physical damage, privacy harm, or public disruption & Operational telemetry, integrity monitoring, physical controls, and anomaly investigation & Layered cyber-physical protection, safe-state design, segmentation, secure maintenance, and rehearsed recovery\\\midrule
Assumptions surrounding physical & Misuse or unexpected state transition & Trust in access is not validated & A local weakness becomes a broader compromise path & Review evidence for threats and test negative paths & Make trust explicit, constrain authority, and fail safely\\\midrule
Recovery capability and decision evidence & Control failure remains unnoticed or cannot be reversed & Monitoring, ownership, or restoration has not been exercised & Longer exposure and slower, less reliable recovery & Alert-to-response exercises and restoration tests & Assigned ownership, actionable telemetry, preserved evidence, and rehearsed recovery\\\midrule
\bottomrule
\end{longtable}
\endgroup
\section*{Practical Security Checklist}
\begin{itemize}[label=$\square$]
\item Define the in-scope assets, users, services, data, and operational dependencies for Physical Security.
\item Document the expected behavior and security objective of physical.
\item Map identities, privileges, trust boundaries, and externally controlled inputs associated with physical and access.
\item Record the threat or failure being addressed, its prerequisites, likely path, and credible consequences.
\item Inspect asset and dependency maps, physical-access records, sensor telemetry, safety constraints, maintenance logs, and recovery exercises.
\item Test both the intended path and negative paths, including invalid state, partial failure, excessive privilege, and unavailable dependencies.
\item Confirm preventive controls are complemented by actionable detection and a named response owner.
\item Exercise containment, restoration, and threats; record actual recovery behavior and unresolved dependencies.
\item Validate exceptions, compensating controls, expiration dates, and accountable approval.
\item Preserve reproducible evidence and tie each finding to affected assets, risk, remediation, and verification criteria.
\end{itemize}
\begin{synthesisbox}
The chapter's principal lesson is that physical security must be evaluated as an operating security system rather than an isolated product or statement. The most important failure patterns involve physical, access, threats, and equipment, especially when ownership, boundaries, telemetry, or recovery remain implicit. A reviewer should connect architecture and behavior to asset and dependency maps, physical-access records, sensor telemetry, safety constraints, maintenance logs, and recovery exercises, prioritize findings by reachable consequence and control independence, and verify remediation through observable change. Within Cyber-Physical Security, this chapter contributes a reusable method: state the objective, expose assumptions, test failure, preserve evidence, and learn from operation.
\end{synthesisbox}
\section*{Self-Test Questions}
\begin{enumerate}
\item What is the central security objective of Physical Security?
\item How does Chapter orientation contribute to the chapter's broader security argument?
\item Distinguish Overview from Physical Security Threats.
\item Which trust assumptions should be made explicit when evaluating physical?
\item How could a weakness involving access affect confidentiality, integrity, availability, privacy, or safety?
\item What evidence would demonstrate that controls surrounding threats operate as intended?
\item Why should prevention, detection, response, and recovery be evaluated together in this domain?
\item Scenario: A connected physical process remains functionally available after a cyber event but no longer operates within safe constraints. What should the reviewer examine first, and why?
\item Scenario: A control associated with equipment passes a configuration check but fails during an operational exercise. How should the finding be interpreted and prioritized?
\item How should a practitioner distinguish enduring principles in this chapter from time-sensitive tools, examples, or implementation details?
\end{enumerate}
\Needspace{8\baselineskip}
\section*{Answer Key}
\begin{enumerate}
\item The objective is to protect the relevant assets by understanding physical consequences, safety, availability, environmental dependencies, remote connectivity, maintenance, and recovery and by connecting controls to measurable risk reduction.
\item Chapter orientation provides one part of the causal chain: it should identify expected behavior, dependencies, failure modes, and the evidence needed to justify confidence.
\item Overview and Physical Security Threats address different portions of the problem. A strong comparison states their scope, assumptions, enforcement points, evidence, and how a failure in one affects the other.
\item Reviewers should identify who controls physical, what inputs or dependencies it trusts, where privilege changes, what happens during partial failure, and how those assumptions are tested.
\item A weakness in access can propagate across trust boundaries. The actual impact depends on reachable assets, granted authority, control independence, visibility, and recovery capability.
\item Useful evidence includes asset and dependency maps, physical-access records, sensor telemetry, safety constraints, maintenance logs, and recovery exercises; configuration alone is insufficient when operation, monitoring, or recovery is part of the claim.
\item Prevention reduces probability, detection limits dwell time, response limits spread, and recovery restores objectives. Omitting one layer creates an unexamined failure mode.
\item Begin by establishing scope, ownership, expected behavior, and the violated assumption. Then evaluate cyber and physical failure together, enforce safe states, and test operational recovery, preserving evidence for prioritization and follow-up.
\item Treat the exercise as stronger evidence than a static check. Prioritize according to reachable impact, likelihood of real operating conditions, missing compensating controls, and recovery difficulty.
\item Retain the principle, threat model, and control objective; separately label technologies, legal references, product examples, and threat observations whose applicability depends on time or environment.
\end{enumerate}
\section*{Key-Term Recap}
\begin{longtable}{>{\RaggedRight\arraybackslash}p{0.27\textwidth}>{\RaggedRight\arraybackslash}p{0.66\textwidth}}
\toprule\rowcolor{Navy}\color{white}\textbf{Term} & \color{white}\textbf{Meaning}\\\midrule
\endfirsthead
\toprule\rowcolor{Navy}\color{white}\textbf{Term} & \color{white}\textbf{Meaning}\\\midrule
\endhead
\textbf{Threat} & A circumstance or actor capable of causing harm by exploiting a weakness or adverse condition. \\\midrule
\textbf{Access Control} & Rules and enforcement mechanisms that restrict actions on resources to authorized subjects. \\\midrule
\textbf{Asset} & Information, service, capability, or physical resource whose value justifies protection. \\\midrule
\textbf{Risk} & The effect of uncertainty on objectives, commonly evaluated through likelihood, impact, exposure, and context. \\\midrule
\textbf{Certificate} & A signed data structure that binds identity information to a public key under a trust model. \\\midrule
\textbf{Policy} & A management statement of required intent, responsibilities, and acceptable behavior. \\\midrule
\textbf{Security Policy} & Documented security intent and requirements that guide decisions and enforcement. \\\midrule
\textbf{Authentication} & The process of establishing confidence that an asserted identity is genuine. \\\midrule
\textbf{Cloud Computing} & On-demand access to pooled computing resources whose security depends on service boundaries and shared responsibilities. \\\midrule
\textbf{Confidentiality} & The property that information is not disclosed to unauthorized parties. \\\midrule
\bottomrule
\end{longtable}
\begin{takeawaybox}
Treat physical security as a verifiable chain from assets and assumptions through controls, evidence, response, and recovery.
\end{takeawaybox}
\end{document}
